\scanpage{219}
\uline{Quelques complexes d'\unsure{Algèbres|algèbres|algèbre} différentielles.}

$A$ anneau commutatif, $M$ $A$-module. \margin{On considère $\mathrm{Sym}^{\ast}(M)\otimes\bigwedge^{\ast}M$ et $\Gamma^{\ast}(M)\otimes\bigwedge^{\ast}M$ \uline{comme} $A$-\uline{alg} \uline{graduées par degré de} $\bigwedge^{\ast}M$}

\margin{De Rham}% slanted
\textcircled{1} Sur $\mathrm{Sym}^{\ast}(M)\otimes\bigwedge^{\ast}M$ \deleted{(algèbre bigraduée)}, il y a un
unique opérateur $d$ qui soit de \deleted{a) $A$-linéaire} \inserted{$A$-algèbre gr. anticommutative, une} \uline{dérivation} \deleted{(pour $\mathrm{Sym}^{\ast}M$ de degré pair et \deleted{$\mathrm{Sym}$} $\bigwedge^{\ast}M$ de degré impair)}
b) qui soit nulle sur $\bigwedge^{1}M$, et qui
sur $\mathrm{Sym}^{1}(M)$ soit l'isom. can. $\mathrm{Sym}^{1}(M)\longrightarrow\bigwedge^{1}M$.

\deleted{\unsure{Il|Le}} Ce $d$ \deleted{\unsure{est}} satisfait $d^2=0$, et applique $\mathrm{Sym}^{p}\otimes\bigwedge^{q}$ dans $\mathrm{Sym}^{p-1}\otimes\bigwedge^{q+1}$. \deleted{Ce $d$ est} \inserted{Il est $\bigwedge^\ast M$-linéaire.}
C'est aussi le complexe de De Rham de $\mathrm{Sym}^{\ast}(M)$ sur $A$, donc le complexe de De Rham du fibré vectoriel $V(M)$ sur $A$.
\deleted{Si \deleted{\unsure{tf}} $M$ est libre et
Si $M$ est libre,}

Supposons $M$ plat\margin{inutile}, \deleted{(par ex., plat tf)} $A$ de car. 0 (i.e. $A$ une $\mathbb Q$-algèbre), alors
\deleted{$H^{\ast}(\mathrm{Sym}^{\ast}(M)\otimes\bigwedge^{\ast}M)$ (gradué par le degré de $\bigwedge^{\ast}M$) est nul en degrés $>0$, et en degré} (graduant suivant le degré de $\bigwedge^\ast M$
\[
H^i_{\mathrm{DR}}(\mathrm{Sym}^{\ast}(M)\otimes\bigwedge^{\ast}(M))\simeq
\begin{cases}0&\text{si }i\ne0\\
\underset{\substack{\shortparallel\\[-2pt]\mathrm{Sym}^{0}\otimes\bigwedge^{0}}}{A}&\text{si }i=0\end{cases}
\]
Si $A$ est de car. $p>0$\inserted{(i.e. est une $\mathbb F_p$-alg), $M$ plat}, alors on a \unsure{.....|....|\ldots} (théorie des opérations de Cartier ----)

\margin{Koszul}% slanted
\textcircled{2} Sur $\mathrm{Sym}^{\ast}(M)\otimes\bigwedge^\ast M$ \deleted{(algèbre bigraduée)} il y a un unique opérateur $\delta$ qui est \deleted{a) $A$-linéaire}
a) une \inserted{anti}dérivation \deleted{si} de \inserted{$A$-alg. graduées, et} \deleted{on prend $\mathrm{Sym}^{1}M$ de degré pair, et b) $\bigwedge^{\ast}M$ de degré impair,}
b) nulle sur \deleted{$\bigwedge^{\unsure{q|0|\ast}}M$} $\mathrm{Sym}^{1}M$, et sur $\bigwedge^{1}M$ c'est l'isom can. $\bigwedge^{1}M\xrightarrow{\sim}\mathrm{Sym}^{1}M$.

Alors $\delta$ satisfait $\delta^2=0$, il applique $\mathrm{Sym}^{p}\otimes\bigwedge^{q}$ dans $\mathrm{Sym}^{p+1}\otimes\bigwedge^{q-1}$. Si $M$ est libre (plus gén, plat) et sans condition de car. sur $A$, on a (\deleted{gr} graduant la coh. par le degré de $\bigwedge^{\ast}$ \inserted{qu'on écrit en bas}
(de sorte que $\delta$ est de degré \deleted{$-1$}\inserted{inf. $1$})):
\[
H^{\mathrm{Kos}}_{\unsure{q|i}}(\mathrm{Sym}^{\ast}M\otimes\bigwedge^\ast M)=
\begin{cases}0&\text{si }i\ne0\\
\simeq\underset{\substack{\shortparallel\\[-2pt]\mathrm{Sym}^{0}\otimes\bigwedge^{0}}}{A}&\text{'' }i=0\end{cases}.
\]
% archive sheet number: 218
